%%%%%%%%%%%%%%%%%%%%%%%%%%%%%%%%%%%%%%%%%%%%%%%%%%%%%%%%%%%%%%%
%% OXFORD MSc THESIS
%%
%% Plain `report' class, styled to match the layout used by recent
%% MSc in Advanced Computer Science submissions:
%%   - 12pt, A4, one-sided, symmetric margins
%%   - one-and-a-half line spacing
%%   - standard "Chapter N / Title" headings (no epigraphs, no minitoc)
%%   - plain centred page numbers, no running heads
%%
%% Everything lives in this one file: preamble, front matter, chapters,
%% and appendices.
%%%%%%%%%%%%%%%%%%%%%%%%%%%%%%%%%%%%%%%%%%%%%%%%%%%%%%%%%%%%%%%

\documentclass[12pt,a4paper]{report}


%%%%% PAGE LAYOUT
% Measured to match the reference layout: 405pt text block, first baseline
% 121pt from the top, last baseline 85pt from the foot, folio 55pt up.
\usepackage[a4paper,hmargin=3.35cm,top=3.83cm,bottom=2.86cm,footskip=26pt]{geometry}

% One-and-a-half spacing for body text.  Captions, footnotes and floats
% stay single-spaced (setspace does this automatically).
\usepackage{setspace}
\onehalfspacing

% Plain centred page number at the foot of every page; no running heads.
\usepackage{fancyhdr}
\pagestyle{fancy}
\fancyhf{}
\renewcommand{\headrulewidth}{0pt}
\fancyfoot[C]{\thepage}
\fancypagestyle{plain}{\fancyhf{}\renewcommand{\headrulewidth}{0pt}\fancyfoot[C]{\thepage}}

% --- raw Unicode symbols used in the draft, mapped to LaTeX (works in text and math) ---
\DeclareUnicodeCharacter{207B}{\ensuremath{^{-}}}
\DeclareUnicodeCharacter{2074}{\ensuremath{^{4}}}
\DeclareUnicodeCharacter{2075}{\ensuremath{^{5}}}
\DeclareUnicodeCharacter{03B2}{\ensuremath{\beta}}
\DeclareUnicodeCharacter{03B5}{\ensuremath{\epsilon}}
\DeclareUnicodeCharacter{03B8}{\ensuremath{\theta}}
\DeclareUnicodeCharacter{03BC}{\ensuremath{\mu}}
\DeclareUnicodeCharacter{03C3}{\ensuremath{\sigma}}
\DeclareUnicodeCharacter{03C4}{\ensuremath{\tau}}
\DeclareUnicodeCharacter{03C6}{\ensuremath{\phi}}
\DeclareUnicodeCharacter{2212}{\ensuremath{-}}
\DeclareUnicodeCharacter{2208}{\ensuremath{\in}}
\DeclareUnicodeCharacter{2248}{\ensuremath{\approx}}
\DeclareUnicodeCharacter{2264}{\ensuremath{\leq}}
\DeclareUnicodeCharacter{2265}{\ensuremath{\geq}}
% --- figure column-image macros (must live in the preamble: definitions inside a figure are local to it) ---
\newcommand{\bimg}[1]{\raisebox{-.5\height}{\includegraphics[width=0.103\linewidth]{sr_baseline_panels/col_#1.png}}}
\newcommand{\bimginp}[1]{\raisebox{-.5\height}{\includegraphics[width=0.103\linewidth]{inp_baseline_panels/col_#1.png}}}
\newcommand{\himg}[2]{\raisebox{-.5\height}{\includegraphics[width=#1\linewidth]{inp_failure/#2.png}}}



%%%%% FONTS AND ENCODING
\usepackage[T1]{fontenc}
\usepackage[utf8]{inputenc}
\usepackage{lmodern}
\usepackage{textcomp}
\usepackage{microtype}
\usepackage[english]{babel}
\usepackage{csquotes}


%%%%% MATHS
\usepackage{amsmath}
\usepackage{amssymb}
\usepackage{amsthm}
\usepackage{bm}

% Uncomment for equation numbers per section (2.3.12) instead of per chapter (2.18):
%\numberwithin{equation}{section}


%%%%% GRAPHICS, TABLES, FLOATS
\usepackage{graphicx}
\graphicspath{{figures/}}
\usepackage{xcolor}
\usepackage{booktabs}
\usepackage{multirow}
\usepackage{array}
\usepackage{longtable}
\usepackage{rotating}
\usepackage{subcaption}
\usepackage[font=small,labelfont=bf]{caption}

% Discourage figure-only pages.
\renewcommand{\topfraction}{0.85}
\renewcommand{\textfraction}{0.1}
\renewcommand{\floatpagefraction}{0.75}
\setlength{\textfloatsep}{20pt plus 15pt minus 4pt}


%%%%% LISTS
\usepackage{enumitem}


%%%%% BIBLIOGRAPHY
% Numeric citations, references listed in order of first appearance.
\usepackage[style=numeric-comp, sorting=none, backend=biber, doi=false, isbn=false]{biblatex}
\addbibresource{references.bib}
\renewcommand*{\bibfont}{\raggedright\small}

\usepackage{notoccite}   % stop citations in captions polluting the bibliography order


%%%%% CROSS-REFERENCES (load hyperref late, cleveref last)
\usepackage[pdfpagelabels]{hyperref}
\usepackage[capitalise,noabbrev]{cleveref}


%%%%% SECTION NUMBERING DEPTH
% 0 = chapter, 1 = section, 2 = subsection, 3 = subsubsection
\setcounter{secnumdepth}{2}   % deepest level that gets a number
\setcounter{tocdepth}{2}      % deepest level shown in the table of contents


%%%%% THESIS DETAILS -- edit these
\newcommand{\thesistitle}{Suitably Impressive Thesis Title}
\newcommand{\candidatenumber}{1234567}
\newcommand{\thesiswordcount}{15000}
\newcommand{\thesisdegree}{Master of Science in Advanced Computer Science}
\newcommand{\thesisdate}{Trinity 2026}


%%%%% PERSONAL MACROS
% A good place to collect your own shorthands as they come up.
% Ordinal superscripts: 4\th place, 1\st, 2\nd, 3\rd
% (\th normally prints the thorn character; we override it, as the Oxford
% template did. Use \textthorn if you ever need the letter.)
\providecommand{\st}{}\renewcommand{\st}{\textsuperscript{st}}
\providecommand{\nd}{}\renewcommand{\nd}{\textsuperscript{nd}}
\providecommand{\rd}{}\renewcommand{\rd}{\textsuperscript{rd}}
\providecommand{\th}{}\renewcommand{\th}{\textsuperscript{th}}


%%%%%%%%%%%%%%%%%%%%%%%%%%%%%%%%%%%%%%%%%%%%%%%%%%%%%%%%%%%%%%%
\begin{document}
%%%%%%%%%%%%%%%%%%%%%%%%%%%%%%%%%%%%%%%%%%%%%%%%%%%%%%%%%%%%%%%

%%%%% TITLE PAGE
\begin{titlepage}
  \centering
  \vspace*{0.14\textheight}

  {\LARGE\bfseries \thesistitle \par}

  \vspace{22mm}
  {\includegraphics[width=3.2cm]{figures/newlogo.pdf} \par}
  \vspace{22mm}

  {\large Candidate Number: \candidatenumber \par}
  \vspace{1.5ex}
  {University of Oxford \par}
  \vspace{1.5ex}
  {Word Count: \thesiswordcount \par}

  \vspace{18mm}

  {A thesis submitted for the degree of \par}
  \vspace{1.5ex}
  {\itshape \thesisdegree \par}
  \vspace{2.5ex}
  {\thesisdate \par}

  \vfill
\end{titlepage}


%%%%%%%%%%%%%%%%%%%%%%%%%%%%%%%%%%%%%%%%%%%%%%%%%%%%%%%%%%%%%%%
%%  FRONT MATTER
%%%%%%%%%%%%%%%%%%%%%%%%%%%%%%%%%%%%%%%%%%%%%%%%%%%%%%%%%%%%%%%
% Roman page numbers, restarting at i.  Page numbers are hidden on the
% abstract and acknowledgements, and become visible from the ToC onwards.
\pagenumbering{roman}
\setcounter{page}{1}

\chapter*{Abstract}
\thispagestyle{empty}

Your abstract goes here. Check your departmental regulations, but this is
usually expected to be under 300 words --- roughly one page.

A workable structure is one paragraph per move: (1) context and why the
problem matters, (2) the gap, (3) what you did, (4) what you found, and
(5) what it means more broadly.

\clearpage


\chapter*{Acknowledgements}
\thispagestyle{empty}

Thank your supervisors first, then anyone whose ideas or feedback shaped
the direction of the work.

Then colleagues, collaborators, and friends.

Then any funding or institutional support you received.

Finally, family.

\clearpage


\tableofcontents
\clearpage

\listoffigures
\clearpage

% Uncomment if you have tables worth listing:
%\listoftables
%\clearpage


%%%%%%%%%%%%%%%%%%%%%%%%%%%%%%%%%%%%%%%%%%%%%%%%%%%%%%%%%%%%%%%
%%  MAIN MATTER
%%%%%%%%%%%%%%%%%%%%%%%%%%%%%%%%%%%%%%%%%%%%%%%%%%%%%%%%%%%%%%%
\clearpage
\pagenumbering{arabic}
\setcounter{page}{1}


\chapter{Introduction}

\section{Motivation}

\section{Contributions}

\section{Thesis Overview}

\chapter{Background}

\section{Generative Modelling}

\subsubsection{The Goal of Generative Modelling}

\subsection{Reward Alignment}

\subsection{Inverse Problems}

\subsubsection{Relation to Reward Alignment}

\section{Flow-based generative models}

\subsection{Normalising Flows}

\subsubsection{Constructing Normalising Flows}

\subsubsection{Training Normalising Flows}

\subsubsection{Forward KL divergence}

\subsubsection{Reverse KL divergence}

\subsection{Flow Matching}

\subsubsection{ODE Definition}

\subsubsection{Flow Model Definition}

\subsubsection{Sampling from a Flow Model}

\subsubsection{Flow Matching}

\subsection{Flow Maps}

\section{Variational Inference}

\subsection{Amortised VI}

\subsection{Stochastic VI}

\section{Variational Flow Maps}

\subsection{Joint training of noise adapter and flow map}

\chapter{Related Work}

\section{The Reward-tilted vector field}

\section{Estimating the Value Function with Flow models}

\subsubsection{Make the reward function robust to noise.}

\subsubsection{Approximating the posterior as the mean}

\subsubsection{Monte Carlo approximation of the posterior}

\subsubsection{Sequential Monte Carlo}

\subsection{Estimating the Value Function with Flow Maps}

\chapter{Method}

\section{The noise adapter as a normalising flow}

\section{Training objectives}

\subsection{Reverse KL}

\subsection{Forward KL}

\subsection{Mixed Objective}

\subsection{$\beta$-VAE}

\subsection{Noise Augmentation}

\section{Building the noise adapter}

\subsection{Low-dimensional Toy Problems}

\subsection{Image-scale Problems}

\section{Evaluation Framework}

\subsubsection{Additional metrics for Low-dimensional experiments}

\subsubsection{Image-scale experiments}

\chapter{Experiments and Results}

\section{2D Checkerboard Example}

\subsection{The Data Distribution}

\subsection{Defining the Inverse Problem}

\subsection{Base Flow Map Training}

\subsection{Noise Adapter Training}

\subsubsection{Normalising Flow Adapter}

\subsubsection{Gaussian Adapter with frozen $\theta$}

\subsubsection{VFM}

\subsection{Results}

\section{Image Inverse Problems}

\subsection{Common Setup}

\subsubsection{Base Flow Map}

\subsubsection{TarFlow Noise Adapter Hyperparameters}

\subsection{Brightness Steering}

\subsubsection{Adapter Training}

\subsubsection{Results}

\subsection{Super-Resolution}

\subsubsection{Noise Adapter Training details}

\subsubsection{Objective Ablation}

\subsubsection{Comparison with Baselines}

\subsection{Inpainting}

\subsubsection{Noise Adapter Training details}

\subsubsection{$\beta$-VAE Ablation}

\subsubsection{Comparison with Baselines}

\subsubsection{Failure Analysis}

\chapter{Conclusion}

\section{Summary}

\section{Future Work}

\section{Critical Evaluation}

\printbibliography[heading=bibintoc,title={Bibliography}]

\appendix

\chapter{Supplementary Material}

\section{Change of Variables}

\section{Score functions and score matching}

\section{Additional Experiments}

\subsection{GMM Experiments}

\subsubsection{Capacity Reward Hack in SR Inverse Problem}

\end{document}
